\clearpage
\section{虚拟体验的心笺、回望与授权}
\label{sec:design-cycle}
心笺和七天回望仅用于虚拟人物体验。用户先记下自己的感受，到回望时再选当前意愿，以及是否授权揭晓。虚拟对方的意愿和授权来自演示情景。

\subsection{从日常记录到当期回望}
聊天过程中可以把感受分数和简短备注记进心笺。七天周期从账号互动空间建立时起算，重置互动状态后重新起算。同一账号共用这个周期，和自然周无关。每到新周期，本人要重新选择意愿和授权，虚拟对方的情景也必须在当期有效。演示时可以推进虚拟日期，观察进入下一期后有哪些状态会变化。

\begin{table}[H]
\centering\small
\caption{虚拟感受记录、回望与揭晓的职责分工}\label{tab:design-cycle-states}
\begin{tabularx}{\linewidth}{@{}p{30mm}X@{}}
\toprule 业务阶段 & 用户操作与处理结果 \\ \midrule
私人心笺 & 自填0--100感受分数及简短备注，保留个人回顾记录 \\
本人回望 & 选择当期意愿及是否授权，为本周期的揭晓提供本人依据 \\
虚拟对方情景 & 提供演示对方的当期意愿与授权，用于虚拟体验中的条件判断 \\
意愿揭晓 & 准入通过后记录本次结果，并使用整个账号本周期的一次机会 \\
\bottomrule
\end{tabularx}
\end{table}

\subsection{授权检查与揭晓结果}
用户发起揭晓时，服务端先检查人物未被屏蔽、双方已有交流，然后检查虚拟对方和本人是否都在当期授权。本人还要选了“想进一步了解”或“有明确好感”，且账号当期机会未使用。若有条件未满足，页面提示缺少的项目。检查通过后，根据演示对方的意愿是否积极，记录双向或非双向结果。

双向和非双向结果均占用当期揭晓机会。比如本期已经对甲揭晓，再想对乙操作就要等到下一期。

